% Generated draft: requires mathematical and source review.
\documentclass[11pt]{article}
\usepackage[margin=25mm]{geometry}
\usepackage{fontspec}
\setmainfont{Latin Modern Roman}
\usepackage{amsmath,amssymb,mathtools,unicode-math}
\setmathfont{Latin Modern Math}
\usepackage{enumitem,needspace,xurl,hyperref}
\setlength{\parindent}{0pt}
\setlength{\parskip}{5pt}
\setlength{\emergencystretch}{3em}
\newcommand{\N}{\mathbb{N}}
\newcommand{\Z}{\mathbb{Z}}
\newcommand{\Q}{\mathbb{Q}}
\newcommand{\R}{\mathbb{R}}
\newcommand{\C}{\mathbb{C}}
\newcommand{\F}{\mathbb{F}}
\newcommand{\A}{\mathbb{A}}
\newcommand{\PP}{\mathbb{P}}
\newcommand{\E}{\mathbb{E}}
\newcommand{\eps}{\varepsilon}
\newcommand{\dd}{\,\mathrm{d}}
\begin{document}
\section*{UnsolvedMath ID: 701695}
\textbf{Source problem code:} catalogue entry
\subsection{Definitions and mathematical statement}
Let \(G\) be a finite group. For subsets \(A,B\subseteq G\), define their product set by
\[
AB=\{xy:x\in A,\ y\in B\}.
\]
Interpret a factorization \(|G|=ab\) as a factorization into positive integers \(a,b\).

The question is whether
\[
\forall\,G\text{ finite group}\quad
\forall\,a,b\in\Z_{\geq 1},\qquad
|G|=ab\ \Longrightarrow\
\exists\,A,B\subseteq G\quad
\bigl(|A|=a\ \land\ |B|=b\ \land\ AB=G\bigr).
\]
The sets \(A\) and \(B\) are arbitrary subsets; no subgroup condition is imposed.

The supplied \emph{Comment of 2025} reports that any minimal counterexample \(G\) must be a simple group without prime-index subgroups, meaning that there is no subgroup \(H\leq G\) for which \([G:H]\) is prime. The source excerpt does not specify the ordering used for minimality, so this contextual assertion is retained with that qualification rather than given a more specific interpretation. The source also refers to Archive problem 19.35, whose statement is not supplied.
\subsection{Short English statement}
Does every factorization of the order of a finite group into two positive integers arise as the sizes of two subsets whose product is the whole group? Proposed by M. H. Hooshmand in 2022. The supplied registry labels the problem open; that status and the accompanying 2025 comment have not been independently verified.
\subsection{Sources}
\begin{itemize}
\item Supplied local UnsolvedMath registry record: permanent ID \(701695\), exact source code \texttt{catalogue entry}, titled \emph{Kourovka Notebook Problem 20.37}. The record supplies the status \texttt{open} and an update timestamp of \texttt{2026-10-16T14:00:00Z}; these are registry metadata, not independent mathematical verification.
\item \emph{The Kourovka Notebook (Unsolved Problems in Group Theory)}, problem 20.37, proposed by M. H. Hooshmand in 2022. The supplied background identifies the 20th issue (2022) and the 21st issue edition (2026), with identifier \texttt{arXiv:1401.0300v48}. The supplied collection description names E. I. Khukhro and V. D. Mazurov as editors.
\item The supplied statement says ``Cf. 19.35 in Archive.'' No statement or further bibliographic information for that archived problem is included.
\item R. R. Bildanov, V. A. Goryachenko, and A. V. Vasil'ev, \emph{Siberian Electron. Math. Rep.}, \textbf{17} (2020), 683--689. Cited by the supplied \emph{Comment of 2025} concerning minimal counterexamples; the article title and a link are not supplied.
\item M. H. Hooshmand, \emph{Commun. Algebra}, \textbf{49}, no. 7 (2021), 2927--2933. Cited by the same supplied \emph{Comment of 2025}; the article title and a link are not supplied.
\end{itemize}
\end{document}
